\begin{frame}[label=q5a3]{5a.3 · DataFrame: región y periodo}
  \small
  \pregunta{¿Qué combinaciones de región y periodo superan las 10\,000 personas reportadas?}
  \entrada{93\,831 filas. \texttt{region} y \texttt{periodo} derivados en el esquema, más \texttt{cantidad}.}
  \flujo{Región y periodo\\Derivados del código}{\texttt{groupBy}\\Dos columnas}{\texttt{SUM(cantidad)}\\Total > 10\,000}{5 combinaciones\\Orden descendente}
  \begin{columns}[T,onlytextwidth]
    \column{0.55\textwidth}
    \textbf{Código Scala original (fragmento)}
    \codigo{sql/05a_3_region_periodo}{11}{17}
    \column{0.42\textwidth}
    \textbf{Salida}\\[5pt]
    {\footnotesize
\begin{tabular}{@{}lll@{}}
\toprule
Región & Periodo & Personas\\
\midrule
Costa & 2022--2024 & 37\,204\\
Costa & 2020--2021 & 24\,603\\
Costa & 2019 & 17\,859\\
Sierra & 2022--2024 & 15\,981\\
\bottomrule
\end{tabular}
    }
  \end{columns}
  \respuesta{Costa 2022--2024 tiene el mayor acumulado. Los periodos duran uno, dos y tres años: comparar estos totales no equivale a comparar promedios anuales.}
  \origen{sql}{05a_3_region_periodo}
\end{frame}
